\documentclass[11pt,a4paper]{article}
\usepackage[margin=1in]{geometry}
\usepackage{amsmath,amssymb,amsthm}
\usepackage{algorithm}
\usepackage{algpseudocode}
\usepackage{booktabs}
\usepackage{hyperref}

\theoremstyle{definition}
\newtheorem{definition}{\textbf{Definition}}[section]
\newtheorem{theorem}{\textbf{Theorem}}[section]
\newtheorem{lemma}[theorem]{\textbf{Lemma}}
\newtheorem{remark}{\textbf{Remark}}[section]

\title{\textbf{Group Relative Policy Optimization (GRPO)}}
\author{\textbf{Team Scaibu} \\ \small Email: \texttt{jaydeep.9664920749@gmail.com} \\ \small \textbf{Architecture and Core Engine Team}}
\date{\textbf{October 2026}}

\begin{document}
\maketitle

\section{Definitions and Cognitive Mental Models}

\subsection{Formal Mathematical Definition}
\textbf{Definition (Critic-Free Group Baselines and Relative Advantage Estimation):}
Standard Proximal Policy Optimization (PPO) in large language model alignment requires maintaining a dedicated value network (critic) $V_\phi(s)$ of comparable scale to the actor policy $\pi_\theta$. In massive parameter regimes (e.g. 70B+ parameters), hosting the critic network consumes equivalent memory and compute to an entire second model, while introducing value estimation instability in complex multi-step reasoning. \textbf{Group Relative Policy Optimization (GRPO)} eliminates the value network entirely by sampling a group of candidate completions for each prompt and computing advantages relative to the group's empirical reward distribution.

For each input prompt $q$, GRPO samples a cohort of $G$ independent completions $\mathcal{O}_q = \{o_1, o_2, \dots, o_G\}$ from the rollout policy $\pi_{\text{old}}$. Each output $o_i$ is evaluated by a reward model or deterministic programmatic verifier, yielding scalar reward $r_i$. The \textbf{Group Relative Advantage} $\hat{A}_i$ is computed without a critic:
{\boldmath
\begin{equation}
\hat{A}_i = \frac{r_i - \bar{r}_q}{\sigma_{r, q} + \epsilon_{\text{std}}}
\end{equation}
}
where $\bar{r}_q = \frac{1}{G}\sum_{i=1}^G r_i$ and $\sigma_{r, q} = \sqrt{\frac{1}{G}\sum_{i=1}^G (r_i - \bar{r}_q)^2}$.

The objective optimizes policy parameters $\theta$ by maximizing the PPO clipped surrogate objective regularized by an unbiased token-level Kullback-Leibler divergence against reference model $\pi_{\text{ref}}$:
{\boldmath
\begin{align}
\mathcal{J}_{\text{GRPO}}(\theta) &= \mathbb{E}_{q \sim P, \{o_i\}_{i=1}^G \sim \pi_{\text{old}}} \left[ \frac{1}{G}\sum_{i=1}^G \left( \frac{1}{|o_i|}\sum_{t=1}^{|o_i|} \min\left( \frac{\pi_\theta(o_{i, t})}{\pi_{\text{old}}(o_{i, t})} \hat{A}_i, \; \operatorname{clip}\left(\frac{\pi_\theta(o_{i, t})}{\pi_{\text{old}}(o_{i, t})}, 1 - \epsilon, 1 + \epsilon\right) \hat{A}_i \right) \right) \right. \nonumber \\
&\quad \left. - \beta D_{\text{KL}}(\pi_\theta \,||\, \pi_{\text{ref}}) \right]
\end{align}
}
where the unbiased KL divergence is computed via the Schulman estimator:
{\boldmath
\begin{equation}
D_{\text{KL}}(\pi_\theta \,||\, \pi_{\text{ref}}) = \frac{\pi_{\text{ref}}(o_{i, t})}{\pi_\theta(o_{i, t})} - \ln \frac{\pi_{\text{ref}}(o_{i, t})}{\pi_\theta(o_{i, t})} - 1
\end{equation}
}

\subsection{Expanded Non-Technical Definition (Intuition for Anyone)}
\textbf{Intuitive Conceptual Definition:}
\begin{enumerate}
    \item \textbf{Fire the Expensive Tutor}: In standard RL (PPO), you have to hire a private full-time tutor (the critic model) whose entire job is to guess how good each word is. The tutor takes up half the computer's memory! GRPO realizes: ``Why hire a tutor? Just let the student try the math problem 8 times, and compare the 8 tries against each other!''
    \item \textbf{Grading on a Curve}: If the student solves a problem 8 times, and gets scores $[10, 10, 0, 0, 0, 0, 0, 0]$, the average is $2.5$. The attempts that scored $10$ receive positive advantages ($+2.5$), while the attempts that failed get negative advantages ($-0.8$). The model naturally learns to repeat the successful attempts.
    \item \textbf{Zero Memory Waste}: By getting rid of the critic model, engineers can use that freed-up GPU memory to make the generator model twice as large, or test 32 reasoning rollouts simultaneously.
    \item \textbf{The Engine of DeepSeek-R1}: GRPO was the key architectural breakthrough created by DeepSeek to train their world-class open-source reasoning models (DeepSeekMath and DeepSeek-R1) to rival OpenAI o1!
\end{enumerate}

\subsection{Child-Friendly Mental Model (Physical Metaphor)}
Imagine an archer practicing at target practice:
\begin{itemize}
    \item \textbf{The Old Method (PPO)}: An expensive coach stands next to the archer, analyzing their elbow angle at every millisecond before the arrow even lands.
    \item \textbf{The GRPO Cohort}: The archer shoots a quiver of 8 arrows at the target.
    \item \textbf{Grading the Quiver}: The archer looks at the target: 2 arrows hit the bullseye, 6 arrows missed to the left.
    \item \textbf{Instant Adjustment}: The archer asks: ``What did my hands feel like on those 2 bullseye shots? Let me do that exact movement again!''
    \item \textbf{Result}: The archer becomes an Olympic champion archer without needing a coach standing by!
\end{itemize}

\section{Core Mathematical Equations and Definitions}

\subsection{Critic-Free Group Advantage Normalization}
{\boldmath
\begin{equation}
\hat{A}_i = \frac{r_i - \frac{1}{G}\sum_{j=1}^G r_j}{\sqrt{\frac{1}{G}\sum_{j=1}^G \left(r_j - \bar{r}\right)^2} + \epsilon_{\text{std}}}
\end{equation}
}
\begin{itemize}
    \item \textbf{Formal Definition}: Z-score standardization of candidate completion rewards within a cohort for prompt $q$.
    \item \textbf{What It Means}: Centers group rewards at zero and scales by standard deviation, giving positive advantage to above-average attempts and negative advantage to below-average attempts.
    \item \textbf{Why It Is Important}: Replaces parameterized state-value estimation with a zero-cost empirical baseline that eliminates value network variance.
\end{itemize}

\subsection{Token-Averaged Group PPO Objective}
{\boldmath
\begin{equation}
\mathcal{L}_{\text{CLIP}}(\theta) = \frac{1}{G}\sum_{i=1}^G \frac{1}{|o_i|}\sum_{t=1}^{|o_i|} \min\left( \frac{\pi_\theta(o_{i, t})}{\pi_{\text{old}}(o_{i, t})} \hat{A}_i, \; \operatorname{clip}\left(\frac{\pi_\theta(o_{i, t})}{\pi_{\text{old}}(o_{i, t})}, 1 - \epsilon, 1 + \epsilon\right) \hat{A}_i \right)
\end{equation}
}
\begin{itemize}
    \item \textbf{Formal Definition}: Clipped surrogate loss averaged first across token positions within each completion, then averaged across group candidates.
    \item \textbf{What It Means}: Normalizes credit assignment across sequences of varying lengths, preventing long rambling chains from dominating gradients.
    \item \textbf{Why It Is Important}: Crucial stability mechanism preventing length explosion in long-chain-of-thought reasoning models.
\end{itemize}

\subsection{Unbiased Token-Level Non-Negative KL Estimator}
{\boldmath
\begin{equation}
D_{\text{KL}}^{\text{unbiased}}(\pi_\theta \,||\, \pi_{\text{ref}}) = \frac{\pi_{\text{ref}}(o_t)}{\pi_\theta(o_t)} - \ln \frac{\pi_{\text{ref}}(o_t)}{\pi_\theta(o_t)} - 1
\end{equation}
}
\begin{itemize}
    \item \textbf{Formal Definition}: Strictly non-negative point-wise divergence estimator proposed by John Schulman.
    \item \textbf{What It Means}: Approximates Kullback-Leibler divergence with zero sample bias and guaranteed non-negativity without requiring full vocabulary expectations.
    \item \textbf{Why It Is Important}: Prevents noisy negative token KL updates from destabilizing the reference constraint during generation.
\end{itemize}

\section{Rigorous Mathematical Analysis Checklist}

\begin{itemize}
    \item \textbf{Notation}: $G$ (cohort size, typically $8 - 64$), $q$ (prompt), $o_i$ (candidate completion), $\hat{A}_i$ (group advantage), $\epsilon$ (clipping threshold).
    \item \textbf{Epistemic Status}: Formulated by Shao et al. (DeepSeek, 2024), core algorithmic engine of DeepSeekMath and DeepSeek-R1.
    \item \textbf{Dependency Graph}: Prompt $q \to \text{Sample } \{o_1, \dots, o_G\} \to \text{Scores } \{r_i\} \to \text{Cohort } \hat{A}_i \to \text{Token Clipping} \to \nabla_\theta \mathcal{J}_{\text{GRPO}}$.
    \item \textbf{Dimensions}: Group size $G$, sequence lengths $|o_i|$.
    \item \textbf{Regimes}: Mathematical reasoning, competitive coding synthesis, formal theorem proving, multi-step agent planning.
    \item \textbf{Invariants}: If all candidate rewards in a group are identical ($r_1 = \dots = r_G$), all advantages $\hat{A}_i \equiv 0.0$ identically and surrogate loss is $0.0$.
    \item \textbf{Isomorphisms}: GRPO is isomorphic to REINFORCE with leave-one-out cohort baselines and importance-sampled PPO trust regions.
\end{itemize}

\section{Theoretical Theorems and Formal Proofs}

\begin{theorem}[Zero-Mean Advantage Invariance Across Cohorts]
For any group size $G \ge 2$ and arbitrary real-valued reward vector $\mathbf{r} = (r_1, \dots, r_G) \in \mathbb{R}^G$, the sum of unnormalized cohort advantages is identically zero:
{\boldmath
\begin{equation}
\sum_{i=1}^G (r_i - \bar{r}) \equiv 0
\end{equation}
}
and consequently, the empirical advantage expectation is strictly unbiased: $\mathbb{E}_{i \sim \operatorname{Unif}(1, G)}[\hat{A}_i] = 0$.
\end{theorem}

\begin{proof}
Let $\bar{r} = \frac{1}{G}\sum_{j=1}^G r_j$.
Summing the advantage deviations across the cohort:
{\boldmath
\begin{align}
\sum_{i=1}^G (r_i - \bar{r}) &= \sum_{i=1}^G r_i - \sum_{i=1}^G \bar{r} \\
&= \sum_{i=1}^G r_i - G \cdot \left( \frac{1}{G}\sum_{j=1}^G r_j \right) \\
&= \sum_{i=1}^G r_i - \sum_{j=1}^G r_j = 0
\end{align}
}
Scaling by the cohort standard deviation $\sigma_r + \epsilon_{\text{std}} > 0$:
{\boldmath
\begin{equation}
\sum_{i=1}^G \hat{A}_i = \frac{1}{\sigma_r + \epsilon_{\text{std}}} \sum_{i=1}^G (r_i - \bar{r}) = 0
\end{equation}
}
Taking the uniform expectation over index $i \in \{1, \dots, G\}$:
{\boldmath
\begin{equation}
\mathbb{E}_i[\hat{A}_i] = \frac{1}{G}\sum_{i=1}^G \hat{A}_i = 0
\end{equation}
}
This proves that group advantage normalization serves as an exact zero-bias baseline for policy gradients, guaranteeing that the baseline neither artificially inflates nor deflates gradient updates.
\end{proof}

\begin{theorem}[Variance Reduction Scaling of Group Baselines]
Let sample rewards $r_i$ be independent and identically distributed with variance $\sigma^2$. The variance of the group baseline estimator $\bar{r}$ scales as $\mathcal{O}\left(\frac{\sigma^2}{G}\right)$, approaching the optimal baseline as $G \to \infty$.
\end{theorem}

\begin{proof}
Let $r_1, \dots, r_G$ be i.i.d. random variables with mean $\mu$ and variance $\sigma^2$.
The group baseline is $\bar{r} = \frac{1}{G}\sum_{i=1}^G r_i$.
Using the standard properties of variance for sums of independent random variables:
{\boldmath
\begin{equation}
\operatorname{Var}(\bar{r}) = \operatorname{Var}\left( \frac{1}{G}\sum_{i=1}^G r_i \right) = \frac{1}{G^2} \sum_{i=1}^G \operatorname{Var}(r_i) = \frac{1}{G^2} (G \sigma^2) = \frac{\sigma^2}{G}
\end{equation}
}
As $G \to \infty$, $\operatorname{Var}(\bar{r}) \to 0$ and $\bar{r} \xrightarrow{p} \mathbb{E}[r(q, \cdot)] = V(q)$.
Thus, for sufficiently large group sizes (e.g. $G \in [16, 64]$), the empirical group mean converges to the true state value $V(q)$, completely eliminating the need for a separate value network approximator.
\end{proof}

\section{Foundational Architectural Questions and Epistemic Meta-Questions}

\subsection{Architectural Question 1: Optimal Group Size Selection}
\textbf{Question:} How is the group size $G$ chosen in production GRPO training?\\
\textbf{Answer:} If $G$ is too small (e.g. $G = 2$), advantage variance is high and cohorts where both attempts fail yield no learning signal. If $G$ is too large (e.g. $G = 256$), generation throughput bottlenecks the cluster. In DeepSeek-R1 and DeepSeekMath, cohorts of $G \in [8, 16]$ for early training and $G \in [32, 64]$ for late reasoning stages achieve an optimal trade-off between baseline precision and accelerator utilization.

\textbf{Meta-Question:} What happens when all completions in a group get the exact same reward?\\
\textbf{Answer:} If all attempts fail ($r_i = 0$) or all attempts succeed ($r_i = 1$), $\sigma_r = 0$ and all advantages $\hat{A}_i = 0$. The prompt contributes zero policy gradient. Filtering prompts by difficulty ensures that training batches prioritize questions with high reward variance (where the model is in its proximal zone of proximal development).

\subsection{Architectural Question 2: Length Normalization Across Reasoning Chains}
\textbf{Question:} Why does GRPO normalize by token length $\frac{1}{|o_i|}$ inside each candidate's sum?\\
\textbf{Answer:} Chain-of-thought solutions for complex math or code vary wildly in length: some solutions take 200 tokens while others take 3,000 tokens. Without token normalization, 3,000-token answers contribute $15\times$ more gradient mass than concise answers, causing the model to generate excessively verbose, hallucinated padding. Token normalization ensures equal gradient weighting across candidate solutions.

\textbf{Meta-Question:} How does GRPO facilitate the emergence of autonomous reasoning behaviors?\\
\textbf{Answer:} Under GRPO with verifiable rewards, the model is rewarded purely for final correctness without human step-by-step demonstrations. Given sufficient compute, the policy autonomously discovers self-reflection, backtracking, and verification strategies (e.g. ``Wait, let me double check my work''), bootstrapping reasoning organically.

\section{Algorithmic Implementation Specification}

\begin{algorithm}[H]
\caption{Group Relative Policy Optimization (GRPO) Evaluation}
\begin{algorithmic}[1]
\Require Rewards $\mathbf{r} \in \mathbb{R}^G$, log-probs $\mathbf{lp}$, old log-probs $\mathbf{lp}_{\text{old}}$, ref log-probs $\mathbf{lp}_{\text{ref}}$, params $\epsilon, \beta$
\Ensure Advantages $\hat{\mathbf{A}}$, mean reward $\bar{r}$, std reward $\sigma_r$, GRPO loss $\mathcal{L}_{\text{GRPO}}$, KL divergence $D_{\text{KL}}$
\State $\bar{r} \gets \frac{1}{G}\sum_{i=0}^{G-1} \mathbf{r}[i], \quad \sigma_r \gets \sqrt{\frac{1}{G}\sum_{i=0}^{G-1} (\mathbf{r}[i] - \bar{r})^2}$
\State Allocate $\hat{\mathbf{A}} \in \mathbb{R}^G$
\For{$i = 0$ \textbf{to} $G-1$}
    \State $\hat{\mathbf{A}}[i] \gets (\mathbf{r}[i] - \bar{r}) / (\sigma_r + 10^{-6})$ \textbf{if} $\sigma_r > 10^{-8}$ \textbf{else} $0.0$
\EndFor
\State $surr\_accum \gets 0.0, \quad kl\_accum \gets 0.0, \quad low \gets 1 - \epsilon, \quad high \gets 1 + \epsilon$
\For{$i = 0$ \textbf{to} $G-1$}
    \State $L_i \gets \operatorname{len}(\mathbf{lp}[i]), \quad seq\_surr \gets 0.0, \quad seq\_kl \gets 0.0$
    \For{$t = 0$ \textbf{to} $L_i - 1$}
        \State $ratio \gets \exp(\mathbf{lp}[i][t] - \mathbf{lp}_{\text{old}}[i][t])$
        \State $clamped \gets \max(low, \min(high, ratio))$
        \State $seq\_surr \gets seq\_surr + \min(ratio \cdot \hat{\mathbf{A}}[i], clamped \cdot \hat{\mathbf{A}}[i])$
        \State $ratio\_ref \gets \exp(\mathbf{lp}_{\text{ref}}[i][t] - \mathbf{lp}[i][t])$
        \State $seq\_kl \gets seq\_kl + \max(0.0, ratio\_ref - (\mathbf{lp}_{\text{ref}}[i][t] - \mathbf{lp}[i][t]) - 1.0)$
    \EndFor
    \State $surr\_accum \gets surr\_accum + seq\_surr / L_i$
    \State $kl\_accum \gets kl\_accum + seq\_kl / L_i$
\EndFor
\State $\mathcal{L}_{\text{GRPO}} \gets - (surr\_accum / G) + \beta \cdot (kl\_accum / G)$
\State \Return $(\hat{\mathbf{A}}, \bar{r}, \sigma_r, \mathcal{L}_{\text{GRPO}}, kl\_accum / G)$
\end{algorithmic}
\end{algorithm}

\section{Big-$\Theta$ Complexity Analysis}

\begin{itemize}
    \item \textbf{Cohort Standardization Time}: $\Theta(G)$ scalar operations across group rewards.
    \item \textbf{Token Surrogate Loss Evaluation}: $\Theta(G \cdot L_{\text{avg}})$ operations across generation tokens.
    \item \textbf{Schulman KL Estimator Time}: $\Theta(G \cdot L_{\text{avg}})$ exponential and subtraction steps.
    \item \textbf{Space Complexity}: $\Theta(G)$ buffer for group advantages.
    \item \textbf{Parallel Reduction Depth}: $\mathcal{D} = \Theta(\log G)$ for cohort advantage normalization.
\end{itemize}

\section{Single Canonical Repository Reference}

The complete deterministic scalar implementation, verification suite, and unit test benchmarks are published at:
\begin{center}
\url{https://github.com/Chief-Strategist-J/llm-obs-infra}
\end{center}

\end{document}
